\documentclass{article}
\usepackage{babel}[spanish]

\title{Examen diagnóstico: Álgebra I}
\begin{document}
\maketitle
\begin{enumerate}
    \item ¿A qué número corresponde la expresión \(5\times 10^6 + 3\times 10^4 + 9\times 10^3 + 2\times 10^2 + 9\times 10^0\)?
    \item Mínimo común múltiplo de 12 y 18
    \item \(5 + 3(8 - 7) - 11(0)\)
    \item \(3x^2\) si \(x = 2\) y \(y = -3\).
    \item \(\frac{5^{-2}\cdot5^4}{5^2}\)
    \item Se compran 25 dulces con \(\mathdollar 12.00\). ¿Cuántos se pueden comprar con \(\mathdollar 36.00\)?
    \item Es la respuesta correcta que resuelve la operación \[\frac{5}{6} - \frac{1}{2} + 2 - \frac{1}{3} = \]
    \item El desarrollo de \((x + 8)(x + 5)\) es \dots
    \item En una tienda se remataron 200 pantalones. Unos de \(\mathdollar 200\) y otros de \(\mathdollar 150\). Si el total de la venta fue de 32500. ¿Cuántos pantalones de 150 se vendieron?
    \item La expresión algebraica del enunciado ``el doble de un número aumentado 3 unidades''
\end{enumerate}
\end{document}
